\section{My impact}

\textbf{3a.} The honest thesis of this section is the line I wrote into our group complaint letter: \emph{the simulation framework we built ourselves is the only reason this project produced any verifiable engineering output at all}. I want to defend that with mechanism, not assertion, because it's exactly the line that at Merit-level reads as self-flattery.

The anchor is a full Kolb (1984) cycle I used deliberately, not retrospectively. My \emph{concrete experience} was three collapsed flight days: the 2026-03-11 weather cancellation, the 2026-03-30 hardware session that produced no airborne minutes, and a third day that ended with our pilot standing on a wet field while Edward and I debugged a corrupt \texttt{calibration\_data.npz}. After the first cancellation my \emph{reflective observation} was harder than ``we lost a day'': the team's entire definition of progress had silently assumed hardware access we didn't have, and the rest of the team was idle because their planned work needed a drone in the air. I was blocked too, but I had a laptop. That asymmetry was the data point. My \emph{abstract conceptualisation} was the move I'm proudest of: if \texttt{vision.py} exposed a single interface (\texttt{detect\_in\_image} returning \texttt{(found, x, y, conf)}) and auto-selected between Ultralytics on laptop and TFLite on Pi, and if \texttt{config.py} auto-detected its hardware, then the same state machine, geofence, and search pattern could run in both places---anything that ran in simulation would, in principle, run on hardware. When I proposed this on the wk18 team call, I wasn't walking away from hardware---I was repositioning the team so that two of us could make verifiable progress without waiting for flight access. Edward kept the Pi bench work alive, Robin built the handoff contracts, Demetro kept the dataset cycle running. The simulator was the thing that let all those tracks produce verifiable output without waiting on each other. My \emph{active experimentation} was six weeks of hammering: \texttt{simple\_simulator.py} grew to 2508 lines with GPS drift and camera shake; \texttt{main.py} was refactored 1411\,$\to$\,754 lines across six modules with 146 pytest tests; and the \texttt{tests/flight/} ladder plus 127 unit tests gave me a way to know, before any field day, which scripts could survive contact with real hardware. When we finally ran on the Pi the delta was thin---a bbox-coordinate bug (commit \texttt{3d62e6a}), a double-scaling fix (\texttt{7eb7cc8}), and a deleted calibration file---and the system came up. I now see the simulation framework wasn't a backup plan but a \emph{statement}: verifiable engineering is possible without hardware access, and simulation-first is no longer a personal preference but a claim about what disciplined engineering owes.

But the sober closing of 3a is that the simulation couldn't tell us what it couldn't tell us. It couldn't predict motion blur at 5\,m/s with a real IMX296---we only modelled Gaussian shake, not directional smear. It couldn't find the descent stall we accepted as a SITL artefact at 4--5\,m. And it couldn't tell us what GPS timing lag (100--200\,ms latency $\to$ 1\,m error at 5\,m/s) would do to a multi-pass lock---that came out of DJI video analysis, after the fact. Sim-to-real parity is an argument I can only win about the things I modelled, and I was disciplined about that only because I'd been forced to be.

Three other contributions matter here and I'll mention them compactly. I unblocked Robin by building him a dedicated HTTP integration path (\texttt{passive\_watch\_2.py}, commit \texttt{a5b1ab7}) with a JSON schema and \texttt{--smart-clear} flag so his state machine could poll vision results without touching my code. I unblocked Demetro by building his two FDR slides from scratch and writing 465 words of speaking script across four iterations. And I drove the group's external communication about hardware failures via the six-iteration complaint letter the whole team signed.

\emph{Against the twelve project requirements R01--R12, our system hit 12/12 in simulation but only 4/12 on real hardware by flight day. Three of the eight gaps were bench-test items we failed to schedule---coordination failure, not engineering failure.} That criterion-judgement-cause chain is the honest self-evaluation I owe: I mistook technical delivery for team enablement, and the 12/12-vs-4/12 split is the proof.

\textbf{3b.} The metric that matters most is one I didn't track until too late: the number of times a teammate other than me edited \texttt{simple\_simulator.py} or \texttt{main.py} was zero. Three concrete mechanisms, each with a test. First, every new tool I build must have at least one named teammate beneficiary documented \emph{before} it ships---if I can't name them, I don't start. Second, no file larger than 500 lines will be refactored solo; I'll pair on the review, and the signal is whether the pair partner can edit that file independently within 48 hours. Third, I'll schedule two hours per week of explicit teaching sessions, graded not by whether I explained something but by whether the teammate completes a handoff task within 48~hours without asking me a second question. If after four weeks those signals aren't firing, the mechanism must change, not my reassurance that I meant well.
